\documentclass[11pt,letterpaper]{article}

\usepackage[margin=1in]{geometry}
\usepackage{graphicx}
\usepackage{booktabs}
\usepackage{amsmath}
\usepackage{hyperref}
\usepackage{xcolor}
\usepackage{caption}
\usepackage{float}

\hypersetup{
  colorlinks=true,
  linkcolor=blue,
  citecolor=blue,
  urlcolor=blue
}

% --- Local macros ---
\newcommand{\pp}{\mbox{$p$+$p$}}
\newcommand{\sqs}{\mbox{$\sqrt{s}$}}
\newcommand{\DR}{\ensuremath{\Delta R}}
\newcommand{\pT}{\ensuremath{p_{\text{T}}}}
\newcommand{\ET}{\ensuremath{E_{\text{T}}}}
\newcommand{\zvtx}{\ensuremath{z_{\mathrm{vtx}}}}
\newcommand{\abseta}{\ensuremath{|\eta|}}
\newcommand{\isoETtruth}{\ensuremath{E_\text{T}^\text{iso,truth}}}
\newcommand{\Rcemc}{\ensuremath{R_{\text{CEMC}}}}
\newcommand{\dzcrit}{\ensuremath{|\Delta z|_{\text{crit}}}}

\graphicspath{{../../plotting/figures/vertex_study/}}

\title{%
  \vspace{-1cm}
  {\large \textbf{sPHENIX Internal Note}}\\[0.5cm]
  {\LARGE \textbf{Vertex Distribution and $\Delta z$ in Single\\
  vs.\ Double-Interaction MC}}\\[0.3cm]
  {\large Isolated Photon Cross-Section in \pp{} at $\sqs = 200$~GeV}
}
\author{sPHENIX Collaboration}
\date{\today}

\begin{document}

\maketitle

\begin{abstract}
This note characterizes the vertex distributions and vertex displacement
($\Delta z = z_{\text{vtx}}^{\text{reco}} - z_{\text{vtx}}^{\text{truth}}$)
in single-interaction and full GEANT double-interaction Monte Carlo
samples used for the sPHENIX isolated photon cross-section measurement.
In double-interaction events, the MBD-reconstructed vertex is the average
of two collision vertices, producing a $\Delta z$ RMS of 48.1~cm (vs.\
16.9~cm for single interaction, before vertex cut). The vertex displacement
maps to a pseudorapidity shift via the EMCal barrel geometry:
$\Delta\eta \approx -\Delta z / \Rcemc$, broadening the cluster
$\Delta\eta$ distribution by a factor of 5 (RMS 0.365 vs.\ 0.073).
We document the vertex distribution shapes, the $\Delta z$ distribution
properties, and the geometric origin of the $\Delta\eta$ correlation
for both interaction types. These results inform the truth-matching
strategy and pileup systematic uncertainty for the cross-section measurement.
\end{abstract}

\tableofcontents
\newpage

%% ============================================================
\section{Introduction}
\label{sec:intro}

In \pp{} collisions at RHIC, the probability of two independent
hard-scattering events occurring in the same bunch crossing (double
interaction, or pileup) depends on the instantaneous luminosity and
crossing angle. For Run~24 at $\sqs = 200$~GeV, the double-interaction
fractions are 18.7\% at 0~mrad crossing angle and 7.2\% at 1.5~mrad.

When two collisions occur in the same bunch crossing, the MBD (Minimum
Bias Detector) reconstructs a single vertex that represents an average
of the two true collision vertices. This displacement between the
reconstructed and true (primary photon) vertex propagates to the
cluster pseudorapidity via the EMCal barrel geometry:
\begin{equation}
  \Delta\eta = \eta_{\text{cluster}} - \eta_{\text{truth}}
    \approx -\frac{\Delta z}{\Rcemc \cdot \cosh\eta}
  \label{eq:deta_dz}
\end{equation}
where $\Rcemc = 93.5$~cm is the EMCal barrel radius.

Understanding the vertex distribution and its consequences is essential
for three aspects of the measurement:
\begin{enumerate}
  \item \textbf{Truth matching:} Vertex displacement drives $\DR$
    matching failures (documented separately in Ref.~\cite{dr_note}).
  \item \textbf{Kinematic migration:} The $\eta$-shift changes both the
    reconstructed $\ET = E / \cosh\eta$ and the cluster's position relative
    to fiducial boundaries, affecting efficiency and the response matrix.
  \item \textbf{Shower shape distortion:} Shifted kinematics alter the
    BDT score evaluation at the wrong working point of the
    $\ET$-dependent threshold.
\end{enumerate}

\begin{thebibliography}{1}
\bibitem{dr_note} sPHENIX Internal Note: Study of the $\DR$
  Truth-Matching Requirement in Single-Interaction MC, 2025.
\end{thebibliography}

%% ============================================================
\section{Samples and Method}
\label{sec:method}

The study uses the macro \texttt{study\_vertex\_double\_interaction.C},
processing two \texttt{photon10} MC samples (truth photon
$14 \leq \pT < 30$~GeV):
\begin{itemize}
  \item \textbf{Single-interaction:}
    \texttt{photon10/bdt\_split.root} --- standard production MC with
    one collision per event.
  \item \textbf{Double-interaction:}
    \texttt{photon10\_double/bdt\_split.root} --- full GEANT4 MC with
    two independent collisions simulated per event.
\end{itemize}

Two million events are processed per sample. Two classes of histograms
are produced:
\begin{enumerate}
  \item \textbf{Event-level} (filled before vertex cut): vertex
    distributions ($z_{\text{vtx}}^{\text{reco}}$,
    $z_{\text{vtx}}^{\text{truth}}$), $\Delta z$ distributions,
    and reco-vs-truth 2D correlation.
  \item \textbf{Per-cluster} (filled after vertex cut and truth
    selection): $\Delta\eta$, $\DR$, $\ET$ ratio, and 2D
    $\Delta\eta$ vs.\ $\Delta z$ correlation for trkID-matched
    truth photons satisfying the standard PPG12 selection
    ($\texttt{pid} = 22$, $\texttt{photonclass} < 3$,
    $\isoETtruth < 4$~GeV, $\abseta < 0.7$,
    $8 \leq \pT < 36$~GeV, $|z_{\text{vtx}}^{\text{reco}}| < 60$~cm).
\end{enumerate}

The distinction matters: event-level distributions include all events
(even those without selected truth photons or outside the vertex cut),
while per-cluster distributions reflect the analysis-relevant population.

%% ============================================================
\section{Results}
\label{sec:results}

%% --------------------------------------------------
\subsection{Reconstructed Vertex Distributions}
\label{sec:vtx_reco}

Figure~\ref{fig:vtx_reco} compares the reconstructed vertex
distributions. The double-interaction reco vertex is \emph{narrower}
than single-interaction, which is a consequence of vertex averaging:
\begin{equation}
  z_{\text{vtx}}^{\text{reco}} \approx \frac{z_1 + z_2}{2}
  \label{eq:vtx_avg}
\end{equation}
Averaging two independent draws from the beam vertex distribution
reduces the variance by $\sqrt{2}$, producing a more peaked reco
distribution.

\begin{figure}[htbp]
\centering
\includegraphics[width=0.75\textwidth]{vertex_reco_comparison.pdf}
\caption{Normalized reconstructed vertex distributions for
  single-interaction (blue) and double-interaction (red) \texttt{photon10}
  MC. The double-interaction distribution is narrower due to the
  averaging of two collision vertices. Both are shown before the
  $|z_{\text{vtx}}| < 60$~cm analysis cut.}
\label{fig:vtx_reco}
\end{figure}

\subsection{Truth Vertex Distributions}
\label{sec:vtx_truth}

The truth vertex (primary collision vertex) distributions are identical
in shape for both samples, as expected --- the beam profile is a
property of the accelerator, not the number of interactions. Both
follow the RHIC luminous region profile at $\sqs = 200$~GeV.

\begin{figure}[htbp]
\centering
\includegraphics[width=0.75\textwidth]{vertex_truth_comparison.pdf}
\caption{Normalized truth vertex distributions (primary collision
  vertex) for single and double-interaction MC. The distributions
  are consistent, reflecting the same RHIC beam profile.}
\label{fig:vtx_truth}
\end{figure}

%% --------------------------------------------------
\subsection{Reco vs.\ Truth Vertex Correlation}
\label{sec:vtx_2d}

Figure~\ref{fig:vtx_2d} shows the 2D correlation between reconstructed
and truth vertex positions. In single-interaction MC, the correlation is
tight along the diagonal with occasional off-diagonal artifacts from MBD
vertex reconstruction failures. In double-interaction MC, the correlation
band is significantly broadened: the reco vertex is pulled toward the
second collision vertex, away from the truth.

\begin{figure}[htbp]
\centering
\includegraphics[width=\textwidth]{vertex_reco_vs_truth.pdf}
\caption{Two-dimensional correlation between reconstructed and truth
  vertex positions for single-interaction (left) and double-interaction
  (right) MC. The red dashed line shows the diagonal
  ($z_{\text{vtx}}^{\text{reco}} = z_{\text{vtx}}^{\text{truth}}$).
  The double-interaction sample shows significant broadening around
  the diagonal due to the second collision vertex.}
\label{fig:vtx_2d}
\end{figure}

%% --------------------------------------------------
\subsection{Vertex Displacement ($\Delta z$) Distributions}
\label{sec:dz}

Figure~\ref{fig:dz} shows the $\Delta z$ distribution for both
samples. The single-interaction distribution has a sharp central
spike (well-reconstructed events) with non-Gaussian tails extending
to $|\Delta z| > 100$~cm from MBD reconstruction failures. The
double-interaction distribution is broad and approximately Gaussian
with heavy tails, reflecting the convolution of two independent
beam vertex distributions.

\begin{table}[htbp]
\centering
\caption{Vertex displacement statistics (event-level, before vertex cut).}
\label{tab:dz_stats}
\begin{tabular}{lcc}
\toprule
\textbf{Metric} & \textbf{Single} & \textbf{Double} \\
\midrule
$\Delta z$ mean & $-0.16$~cm & $+0.13$~cm \\
$\Delta z$ RMS  & 16.9~cm    & 48.1~cm \\
Fraction $|\Delta z| > 9.35$~cm & 36.0\% & 74.2\% \\
\bottomrule
\end{tabular}
\end{table}

The critical threshold $\dzcrit = \DR_{\text{cut}} \times \Rcemc = 0.1
\times 93.5 = 9.35$~cm marks where the vertex-induced $\Delta\eta$
exceeds the $\DR = 0.1$ matching cut (at $\eta = 0$). In
double-interaction MC, 74.2\% of events exceed this threshold, compared
to 36.0\% in single-interaction (before vertex cut). The large
single-interaction fraction is dominated by events outside the
$|z_{\text{vtx}}| < 60$~cm cut; within the analysis acceptance, the
effective $\DR$ failure rate drops to 3.04\%.

\begin{figure}[htbp]
\centering
\includegraphics[width=0.75\textwidth]{dz_distribution_comparison.pdf}
\caption{Normalized $\Delta z$ distributions for single (blue) and
  double (red) interaction MC. Dashed lines mark the critical threshold
  $|\Delta z| = 9.35$~cm. The double-interaction distribution is
  dramatically broader, with most events exceeding the threshold.}
\label{fig:dz}
\end{figure}

\begin{figure}[htbp]
\centering
\includegraphics[width=0.75\textwidth]{dz_abs_comparison.pdf}
\caption{Normalized $|\Delta z|$ distributions. The single-interaction
  distribution drops steeply beyond 5~cm, while the double-interaction
  distribution remains approximately flat out to $\sim$60~cm. The
  critical threshold $\dzcrit = 9.35$~cm is marked.}
\label{fig:dz_abs}
\end{figure}

%% --------------------------------------------------
\subsection{Physics of the $\Delta z$ Distribution}
\label{sec:dz_physics}

For double-interaction events, if both collision vertices are drawn
independently from the beam luminous region with distribution $f(z)$
(approximately Gaussian, $\sigma_z \approx 30$~cm for RHIC at
$\sqs = 200$~GeV), then:
\begin{equation}
  \Delta z = z_{\text{vtx}}^{\text{reco}} - z_1
    \approx \frac{z_1 + z_2}{2} - z_1 = \frac{z_2 - z_1}{2}
  \label{eq:dz_double}
\end{equation}

The distribution of $(z_2 - z_1)$ for two independent draws from $f(z)$
has variance $2\sigma_z^2$, so:
\begin{equation}
  \sigma_{\Delta z}^{\text{double}} = \frac{\sqrt{2}\,\sigma_z}{2}
    = \frac{\sigma_z}{\sqrt{2}}
  \label{eq:dz_sigma}
\end{equation}

With $\sigma_z \approx 30$~cm, the predicted $\sigma_{\Delta z}
\approx 21$~cm. The measured value of 48.1~cm is larger, reflecting
the non-Gaussian tails of the RHIC luminous region and the contribution
of events where the MBD vertex reconstruction is biased beyond simple
averaging.

For single-interaction events, $\Delta z$ arises solely from MBD
vertex reconstruction resolution. The sharp central peak
($\sigma \sim 3$~cm) represents well-reconstructed events, while the
tails reflect events with insufficient MBD timing information for
precise vertex determination.

%% --------------------------------------------------
\subsection{$\Delta\eta$ vs.\ $\Delta z$ Correlation}
\label{sec:deta_dz}

Figure~\ref{fig:deta_dz} shows the 2D correlation between cluster
$\Delta\eta$ and vertex displacement $\Delta z$ for double-interaction
MC. The tight linear correlation confirms the geometric origin
(Eq.~\ref{eq:deta_dz}), with the finite band width arising from the
$\cosh\eta$ factor varying across the barrel ($|\eta| < 0.7$).

\begin{figure}[htbp]
\centering
\includegraphics[width=0.75\textwidth]{deta_vs_dz_double.pdf}
\caption{Two-dimensional $\Delta\eta$ vs.\ $\Delta z$ for
  double-interaction MC (\texttt{photon10\_double}, trkID-matched
  truth photons). The red dashed line shows
  $\Delta\eta = -\Delta z / \Rcemc$ at $\eta = 0$. Dotted lines mark
  $\Delta\eta = \pm 0.1$, corresponding to the $\DR$ matching threshold
  at $\Delta\phi \approx 0$. The band width reflects the $\cosh\eta$
  variation across the barrel acceptance.}
\label{fig:deta_dz}
\end{figure}

%% --------------------------------------------------
\subsection{$\Delta\eta$ Distribution}
\label{sec:deta}

Figure~\ref{fig:deta} compares the $\Delta\eta$ distributions. The
single-interaction distribution is narrowly peaked (RMS 0.073), while
the double-interaction distribution extends to $|\Delta\eta| > 1.0$
(RMS 0.365, a factor of 5 broader). This broadening directly propagates
to $\ET$ migration and BDT threshold evaluation.

\begin{figure}[htbp]
\centering
\includegraphics[width=0.75\textwidth]{deta_comparison.pdf}
\caption{Normalized $\Delta\eta$ distributions for single (blue) and
  double (red) interaction MC. The double-interaction distribution
  extends well beyond $|\Delta\eta| = 0.1$ (the approximate $\DR$
  threshold), with RMS 5$\times$ larger.}
\label{fig:deta}
\end{figure}

%% --------------------------------------------------
\subsection{Impact on $\DR$ and $\ET$ Matching}
\label{sec:impact}

Table~\ref{tab:matching} summarizes the matching statistics for
truth photons within the analysis acceptance
($|z_{\text{vtx}}^{\text{reco}}| < 60$~cm).

\begin{table}[htbp]
\centering
\caption{Truth-matching statistics (within vertex cut, standard
  PPG12 truth photon selection).}
\label{tab:matching}
\begin{tabular}{lcc}
\toprule
\textbf{Metric} & \textbf{Single} & \textbf{Double} \\
\midrule
Truth photons selected  & 476{,}381 & 690{,}450 \\
trkID matched           & 465{,}270 (97.7\%) & 635{,}532 (92.0\%) \\
Fail $\DR < 0.1$        & 14{,}164 (3.04\%) & 413{,}620 (65.08\%) \\
Cluster $\Delta\eta$ RMS & 0.073 & 0.365 \\
\bottomrule
\end{tabular}
\end{table}

\begin{figure}[htbp]
\centering
\includegraphics[width=0.75\textwidth]{deltaR_comparison.pdf}
\caption{Normalized $\DR$ distributions for trkID-matched clusters
  in single (blue) and double (red) interaction MC. The dashed line
  marks the $\DR = 0.1$ cut. In double-interaction MC, the majority
  of correctly trkID-matched clusters exceed the geometric threshold.}
\label{fig:deltaR}
\end{figure}

\begin{figure}[htbp]
\centering
\includegraphics[width=0.75\textwidth]{ET_ratio_comparison.pdf}
\caption{Distribution of $\ET^{\text{cluster}} / \pT^{\text{truth}}$
  for trkID-matched clusters. The double-interaction distribution
  is broadened by the $\ET$ shift from the displaced vertex,
  but remains centered near unity --- indicating the calorimeter
  correctly measured the photon energy.}
\label{fig:ET_ratio}
\end{figure}

The $\ET$ ratio distribution (Figure~\ref{fig:ET_ratio}) shows that
despite the large vertex displacement, the cluster energy measurement
remains correct. The broadening of the $\ET / \pT^{\text{truth}}$
distribution in double-interaction MC reflects the $\ET$ recalculation
at the wrong $\eta$ (due to the displaced vertex), not a degradation
of calorimeter response.

%% --------------------------------------------------
\subsection{$\Delta\eta$ Dependence on $\eta$}
\label{sec:deta_eta}

The $\Delta\eta$ shift depends on the cluster position in the barrel
through the $\cosh\eta$ factor in Eq.~(\ref{eq:deta_dz}). At
$\eta = 0$, $\cosh\eta = 1$ and the shift is maximal for a given
$\Delta z$. At $|\eta| = 0.7$, $\cosh\eta \approx 1.26$, reducing
the shift by 20\%.

Figure~\ref{fig:deta_eta} shows the $\Delta\eta$ vs.\ $\eta_{\text{truth}}$
correlation for both samples. In single-interaction MC, the correlation
is barely visible. In double-interaction MC, the characteristic ``X''
pattern reflects the $\cosh\eta$ modulation.

\begin{figure}[htbp]
\centering
\includegraphics[width=\textwidth]{deta_vs_eta_comparison.pdf}
\caption{$\Delta\eta$ vs.\ truth $\eta$ for single (left) and double
  (right) interaction MC. Note the different $\Delta\eta$ scales. In
  double interaction, the spread is modulated by $\cosh\eta$, with
  maximum displacement at $\eta = 0$.}
\label{fig:deta_eta}
\end{figure}

%% ============================================================
\section{Summary}
\label{sec:summary}

\begin{enumerate}
  \item The double-interaction MBD vertex is the average of two
    collision vertices, producing a $\Delta z$ RMS of 48.1~cm (vs.\
    16.9~cm for single interaction, event-level before vertex cut).

  \item The reconstructed vertex distribution is actually \emph{narrower}
    for double-interaction events (due to averaging), but the
    displacement from the truth vertex is much larger.

  \item The $\Delta\eta$ shift follows the geometric prediction
    $\Delta\eta \approx -\Delta z / (\Rcemc \cdot \cosh\eta)$,
    with cluster $\Delta\eta$ RMS broadened by a factor of 5
    (0.365 vs.\ 0.073).

  \item 74.2\% of double-interaction events exceed the critical
    vertex displacement $|\Delta z| = 9.35$~cm, leading to a
    65.1\% $\DR$-matching failure rate for trkID-matched clusters.

  \item Cluster energy response ($\ET / \pT^{\text{truth}}$) remains
    centered near unity for both interaction types; the vertex shift
    distorts the angular position, not the energy measurement.

  \item These results provide the quantitative basis for:
    (a) removing the geometric $\DR$ matching cut in favor of
    trkID-only truth association;
    (b) evaluating the pileup systematic through the double-interaction
    efficiency ratio; and
    (c) understanding the vertex-driven component of shower-shape
    BDT score migration.
\end{enumerate}

%% ============================================================
\section*{Files}
\label{sec:files}

\begin{tabular}{ll}
\toprule
\textbf{File} & \textbf{Purpose} \\
\midrule
\texttt{efficiencytool/study\_vertex\_double\_interaction.C} & Study macro \\
\texttt{plotting/plot\_vertex\_double\_interaction.C} & Plotting macro \\
\texttt{results/vertex\_double\_interaction\_study.root} & Output histograms \\
\texttt{plotting/figures/vertex\_study/*.pdf} & Output figures \\
\bottomrule
\end{tabular}

\end{document}
